% Source/coverage: Weinberg QFT Vol. I, Chapter 8 Appendix, physical PDF pp. 395--397
% (printed pp. 372--374), from "Appendix Traces" through the sentence before Problems.
% Equations (8.A.1)--(8.A.12); one ordinary source note; no Problems transcribed.

\chapterappendix{}{Traces}

In calculating $S$-matrix elements and transition rates for processes
involving particles of spin $\tfrac12$, we often encounter traces of products
of Dirac gamma matrices. It will therefore be useful to give formulas for
these traces that can be used in all such calculations.

For products of even numbers of gamma matrices, the trace is given by
\begin{equation}
\operatorname{Tr}\{\gamma_{\mu_1}\gamma_{\mu_2}\cdots\gamma_{\mu_{2N}}\}
=4\sum_{\text{pairings}}\delta_P
\prod_{\text{pairs}}\eta_{\text{paired }\mu\text{s}}.
\tag{8.A.1}\label{eq:8.A.1}
\end{equation}
Here the sum is over all different ways of pairing the indices
$\mu_1,\ldots,\mu_{2N}$. A pairing can be regarded as a permutation of the
integers $1,2,\ldots,2N$ into some order $P1,P2,\ldots,P(2N)$,
in which we pair $\mu_{P1}$ with $\mu_{P2}$, $\mu_{P3}$ with $\mu_{P4}$, and
so on. Permuting pairs or permuting $\mu$s within a pair yields the same
pairing, so the number of different pairings is
\begin{equation}
\frac{(2N)!}{N!\,2^N}=(2N-1)(2N-3)\cdots1\equiv(2N-1)!!.
\tag{8.A.2}\label{eq:8.A.2}
\end{equation}
We can avoid summing over equivalent pairings by requiring that
\begin{equation}
P1<P2,\quad P3<P4,\quad\ldots,\quad
P(2N-1)<P(2N)
\tag{8.A.3}\label{eq:8.A.3}
\end{equation}
and
\begin{equation}
P1<P3<P5<\cdots.
\tag{8.A.4}\label{eq:8.A.4}
\end{equation}
With this convention, the factor $\delta_P$ is $+1$ or $-1$ according to
whether the pairing involves an even or odd permutation of indices. The
product in Eq.~\eqref{eq:8.A.1} is over all $N$ pairs, the $n$th pair contributing a
factor $\eta_{\mu_{P(2n-1)}\mu_{P(2n)}}$. For instance (writing
$\mu,\nu,\rho,\sigma,\ldots$ in place of
$\mu_1,\mu_2,\mu_3,\mu_4,\ldots$), for $N=1$, 2, and 3 we have\footnote{There
are now computer programs~\hyperref[ref:8.7]{[7]} available for the calculation of traces of
products of large numbers of Dirac matrices.}
\begin{equation}
\operatorname{Tr}\{\gamma_\mu\gamma_\nu\}=4\eta_{\mu\nu},
\tag{8.A.5}\label{eq:8.A.5}
\end{equation}
\begin{equation}
\operatorname{Tr}\{\gamma_\mu\gamma_\nu\gamma_\rho\gamma_\sigma\}
=4[\eta_{\mu\nu}\eta_{\rho\sigma}
-\eta_{\mu\rho}\eta_{\nu\sigma}
+\eta_{\mu\sigma}\eta_{\nu\rho}],
\tag{8.A.6}\label{eq:8.A.6}
\end{equation}
\begin{align}
&\operatorname{Tr}\{\gamma_\mu\gamma_\nu\gamma_\rho\gamma_\sigma
\gamma_\kappa\gamma_\eta\}=4\bigl[
\eta_{\mu\nu}\eta_{\rho\sigma}\eta_{\kappa\eta}
-\eta_{\mu\nu}\eta_{\rho\kappa}\eta_{\sigma\eta}
+\eta_{\mu\nu}\eta_{\rho\eta}\eta_{\sigma\kappa}
\notag\\
&\quad-\eta_{\mu\rho}\eta_{\nu\sigma}\eta_{\kappa\eta}
+\eta_{\mu\rho}\eta_{\nu\kappa}\eta_{\sigma\eta}
-\eta_{\mu\rho}\eta_{\nu\eta}\eta_{\sigma\kappa}
+\eta_{\mu\sigma}\eta_{\nu\rho}\eta_{\kappa\eta}
-\eta_{\mu\sigma}\eta_{\nu\kappa}\eta_{\rho\eta}
\notag\\
&\quad+\eta_{\mu\sigma}\eta_{\nu\eta}\eta_{\rho\kappa}
-\eta_{\mu\kappa}\eta_{\nu\rho}\eta_{\sigma\eta}
+\eta_{\mu\kappa}\eta_{\nu\sigma}\eta_{\rho\eta}
-\eta_{\mu\kappa}\eta_{\nu\eta}\eta_{\rho\sigma}
+\eta_{\mu\eta}\eta_{\nu\rho}\eta_{\sigma\kappa}
\notag\\
&\quad-\eta_{\mu\eta}\eta_{\nu\sigma}\eta_{\rho\kappa}
+\eta_{\mu\eta}\eta_{\nu\kappa}\eta_{\rho\sigma}\bigr].
\tag{8.A.7}\label{eq:8.A.7}
\end{align}
For an odd number of gamma matrices, the result is much simpler
\begin{equation}
\operatorname{Tr}\{\gamma_{\mu_1}\gamma_{\mu_2}\cdots
\gamma_{\mu_{2N+1}}\}=0.
\tag{8.A.8}\label{eq:8.A.8}
\end{equation}

The proof of Eq.~\eqref{eq:8.A.1} is by mathematical induction. First note that
\[
\operatorname{Tr}\{\gamma_\mu\gamma_\nu\}
=-\operatorname{Tr}\{\gamma_\nu\gamma_\mu\}
+2\operatorname{Tr}\{\eta_{\mu\nu}\mathbf{1}\}
=-\operatorname{Tr}\{\gamma_\mu\gamma_\nu\}+8\eta_{\mu\nu},
\]
so $\operatorname{Tr}\{\gamma_\mu\gamma_\nu\}=4\eta_{\mu\nu}$, in agreement
with Eq.~\eqref{eq:8.A.1}. Next, suppose that Eq.~\eqref{eq:8.A.1} is true for $N\leq M-1$.
We then have
\begin{align*}
\operatorname{Tr}\{\gamma_{\mu_1}\gamma_{\mu_2}\cdots\gamma_{\mu_{2M}}\}
={}&2\eta_{\mu_1\mu_2}
\operatorname{Tr}\{\gamma_{\mu_3}\cdots\gamma_{\mu_{2M}}\}
-\operatorname{Tr}\{\gamma_{\mu_2}\gamma_{\mu_1}
\gamma_{\mu_3}\cdots\gamma_{\mu_{2M}}\}\\
={}&2\eta_{\mu_1\mu_2}
\operatorname{Tr}\{\gamma_{\mu_3}\cdots\gamma_{\mu_{2M}}\}
-2\eta_{\mu_1\mu_3}
\operatorname{Tr}\{\gamma_{\mu_2}\gamma_{\mu_4}\cdots\gamma_{\mu_{2M}}\}\\
&+\operatorname{Tr}\{\gamma_{\mu_2}\gamma_{\mu_3}\gamma_{\mu_1}
\gamma_{\mu_4}\cdots\gamma_{\mu_{2M}}\}\\
={}&2\eta_{\mu_1\mu_2}
\operatorname{Tr}\{\gamma_{\mu_3}\cdots\gamma_{\mu_{2M}}\}
-2\eta_{\mu_1\mu_3}
\operatorname{Tr}\{\gamma_{\mu_2}\gamma_{\mu_4}\cdots\gamma_{\mu_{2M}}\}\\
&+2\eta_{\mu_1\mu_4}
\operatorname{Tr}\{\gamma_{\mu_2}\gamma_{\mu_3}\gamma_{\mu_5}
\cdots\gamma_{\mu_{2M}}\}-\cdots\\
&+2\eta_{\mu_1\mu_{2M}}
\operatorname{Tr}\{\gamma_{\mu_2}\cdots\gamma_{\mu_{2M-1}}\}
-\operatorname{Tr}\{\gamma_{\mu_2}\cdots\gamma_{\mu_{2M}}\gamma_{\mu_1}\}.
\end{align*}
All commutators have zero trace, so the last term subtracted here is the same
as the left-hand side, and so
\begin{align}
\operatorname{Tr}\{\gamma_{\mu_1}\gamma_{\mu_2}\cdots\gamma_{\mu_{2M}}\}
={}&\eta_{\mu_1\mu_2}
\operatorname{Tr}\{\gamma_{\mu_3}\cdots\gamma_{\mu_{2M}}\}
-\eta_{\mu_1\mu_3}
\operatorname{Tr}\{\gamma_{\mu_2}\gamma_{\mu_4}\cdots\gamma_{\mu_{2M}}\}
\notag\\
&+\eta_{\mu_1\mu_4}
\operatorname{Tr}\{\gamma_{\mu_2}\gamma_{\mu_3}\gamma_{\mu_5}
\cdots\gamma_{\mu_{2M}}\}-\cdots
\notag\\
&+\eta_{\mu_1\mu_{2M}}
\operatorname{Tr}\{\gamma_{\mu_2}\cdots\gamma_{\mu_{2M-1}}\}.
\tag{8.A.9}\label{eq:8.A.9}
\end{align}
If we assume that Eq.~\eqref{eq:8.A.1} correctly gives the trace of any product of
$2N-2$ Dirac matrices, then Eq.~\eqref{eq:8.A.9} shows that Eq.~\eqref{eq:8.A.1} also
correctly gives the trace of any product of $2N$ Dirac matrices.

The easiest way to see that the trace of an odd number of Dirac matrices
vanishes is to note that $-\gamma_\mu$ is related to $\gamma_\mu$ by a
similarity transformation,
$-\gamma_\mu=\gamma_5\gamma_\mu(\gamma_5)^{-1}$. Traces are unaffected by
such similarity transformations, so the trace of an odd number of Dirac
matrices is equal to minus itself, and hence vanishes.

One occasionally encounters another class of traces, of the form
\[
\operatorname{Tr}\{\gamma_5\gamma_{\mu_1}\gamma_{\mu_2}\cdots
\gamma_{\mu_n}\}.
\]
This vanishes for odd $n$ for the same reason as given above for traces
without a $\gamma_5$. It also vanishes for $n=0$ and $n=2$:
\begin{equation}
\operatorname{Tr}\{\gamma_5\}=0
\tag{8.A.10}\label{eq:8.A.10}
\end{equation}
\begin{equation}
\operatorname{Tr}\{\gamma_5\gamma_\mu\gamma_\nu\}=0.
\tag{8.A.11}\label{eq:8.A.11}
\end{equation}
(To see this just recall that
$\gamma_5\equiv \ii\gamma^0\gamma^1\gamma^2\gamma^3$, and note that there is no
way of pairing the indices in
$\operatorname{Tr}\{\gamma_0\gamma_1\gamma_2\gamma_3\}$ or in
$\operatorname{Tr}\{\gamma_0\gamma_1\gamma_2\gamma_3\gamma_\mu\gamma_\nu\}$
so that the spacetime indices in each pair are equal.) For $n=4$ it is
possible to pair the indices in
$\operatorname{Tr}\{\gamma_0\gamma_1\gamma_2\gamma_3
\gamma_\mu\gamma_\nu\gamma_\rho\gamma_\sigma\}$ so that the spacetime
indices in each pair are equal, but only if $\mu,\nu,\rho,\sigma$ are some
permutation of $0,1,2,3$. Furthermore this trace must be odd under
permutations of $\mu,\nu,\rho,\sigma$ since gamma matrices with different
indices anticommute. Thus the trace
$\operatorname{Tr}\{\gamma_5\gamma_\mu\gamma_\nu\gamma_\rho\gamma_\sigma\}$
must be proportional to the totally antisymmetric tensor
$\epsilon_{\mu\nu\rho\sigma}$. The constant of proportionality may be worked
out by letting $\mu,\nu,\rho,\sigma$ take the values $0,1,2,3$, and recalling
that $\epsilon_{0123}=-1$. In this way we find
\begin{equation}
\operatorname{Tr}\{\gamma_5\gamma_\mu\gamma_\nu\gamma_\rho\gamma_\sigma\}
=4\ii\epsilon_{\mu\nu\rho\sigma}.
\tag{8.A.12}\label{eq:8.A.12}
\end{equation}
The trace of products of $\gamma_5$ with six, eight, or more Dirac matrices
may be calculated by the same methods used above to verify Eq.~\eqref{eq:8.A.1}.
